% !PRINTED 515
\weberpage{515}

Denn nach der Voraussetzung über $\alpha, \beta$ ist
\[ \alpha x + \beta y = \varepsilon \]
eine Einheit. Durch Multiplication mit $\mu$ folgt daraus
\[ \alpha \mu x + \beta \mu y = \varepsilon \mu, \]
und da $\alpha \mu$ und $\beta \mu$ nach Voraussetzung durch $\beta$ theilbar sind, so ist auch $\varepsilon \mu$ und folglich auch das mit $\varepsilon \mu$ associirte $\mu$ durch $\beta$ theilbar.

\begin{center}
\S.\,140.

\gesperrt{Primfunctionale im Körper $\Omega$.}
\end{center}

Durch die Sätze des vorigen Paragraphen haben wir die Hülfsmittel gewonnen, um die Gesetze der Theilbarkeit der ganzen Functionale im Körper $\Omega$ genau auf demselben Wege abzuleiten, den man in den Elementen der Arithmetik auf die natürlichen ganzen Zahlen anwendet. Wir definiren folgendermaassen:

1. \gesperrt{Ein ganzes Functional $\pi$ des Körpers $\Omega$, welches keine Einheit ist, heisst ein Primfunctional, wenn es ausser durch die Einheiten nur noch durch die mit ihm selbst associirten Functionale theilbar ist. Jedes ganze Functional in $\Omega$, das ausser diesen noch andere Theiler hat, heisst zusammengesetzt.}

Der Begriff des Primfunctionales ist hiernach wesentlich von dem Körper $\Omega$ abhängig. Es können sehr wohl die Primfunctionale eines Körpers in einem anderen erweiterten Körper zusammengesetzt sein.

Wenn $\omega$ ein beliebiges ganzes und $\pi$ ein Primfunctional des Körpers $\Omega$ ist, so sind nur zwei Fälle möglich: entweder $\omega$ ist relativ prim zu $\pi$ oder $\omega$ ist durch $\pi$ theilbar; denn ein gemeinschaftlicher Theiler von $\omega$ und $\pi$ kann nur entweder eine Einheit oder mit $\pi$ associirt sein, und im letzteren Falle ist $\omega$ durch $\pi$ theilbar. Daraus ergiebt sich der Satz:

2. \gesperrt{Wenn das Product $\alpha \beta$ zweier ganzer Functionale $\alpha, \beta$ in $\Omega$ durch ein Primfunctional $\pi$ theilbar ist, so muss einer der beiden Factoren durch $\pi$ theilbar sein.}
